% !TEX root = ../main.tex
\chapter{関係・前順序・Galois 接続}
\label{chap:relations-preorders-galois}

\begin{chapterabstract}
本章では、等式だけでは扱いきれない仕様の比較を扱う。ソフトウェアでは、二つの実装やデータが完全に同じでなくても、一方が他方より詳しい、より制約が強い、より安全である、より抽象的である、と言える場面が多い。このような比較を表す道具が関係と前順序である。さらに、詳細な世界と抽象的な世界を往復させるとき、抽象化と具体化の対応を Galois 接続として整理できる。本章では、公開ログへのマスキングを題材に、情報を落とす処理を「失敗した同型」ではなく「安全な近似」として読む方法を学ぶ。
\end{chapterabstract}

\begin{learninggoals}
\item 等式ではなく関係で仕様を表す場面を説明できる。
\item 前順序を、反射性と推移性を持つ比較として読める。
\item 仕様の強弱、情報量、権限、抽象度を順序としてモデル化できる。
\item 単調性を「比較を壊さない変換」として理解できる。
\item Galois 接続を、抽象化と具体化がずれなく対応するための法則として説明できる。
\end{learninggoals}

\section{等式では足りない場面}

第12章から第15章までは、主に「同じ結果になること」を扱ってきた。リファクタリング前後の式が同じであること、関数がすべての入力で同じ出力を返すこと、adapter が業務ロジックと干渉しないことは、等式で表しやすい。

しかし実務では、等式だけでは言い切れない安全性が多い。たとえば、詳細ログから公開ログを作る処理を考える。公開ログは、詳細ログと同じではない。ユーザーID、内部トークン、IPアドレス、細かなエラー理由などを落とすかもしれない。したがって、詳細ログと公開ログの間に同型を期待してはいけない。

それでも、公開ログへの変換は正しいと言いたい場合がある。言いたいことは「元に戻せる」ではない。「公開してはいけない情報を含まない」「公開側に残すと決めた項目は保存されている」「詳細なログを抽象化したものとして妥当である」といった主張である。これは等式ではなく、関係や順序の言葉で表すほうが自然である。

\begin{intuitionbox}
等式は「同じ」を表す。関係は「何らかの意味で対応している」を表す。前順序は「一方が他方以下である」「一方が他方より弱い」「一方が他方より抽象的である」といった片方向の比較を表す。
\end{intuitionbox}

\begin{softwarebox}
情報を落とす変換を、無理に round-trip として扱わない。公開ログ、権限、静的解析、lossy migration では、「戻せるか」ではなく「安全な方向に近似しているか」を確認する。
\end{softwarebox}

\FigurePlaceholder{fig-ch16-equality-vs-order.png}{0.90}{等式で表す意味保存と、順序で表す安全な近似の違い}{fig:ch16-equality-vs-order}

\section{関係：二つの値を比べる命題}

最初の道具は関係である。関係とは、二つの値を受け取って命題を返すものだと考えればよい。

たとえば、詳細ログ \texttt{DetailLog} と公開ログ \texttt{PublicLog} の間に、「公開ログの endpoint は詳細ログの endpoint と一致している」という関係を定義できる。これは二つのログが等しいという主張ではない。公開ログが詳細ログの一部を正しく反映しているという主張である。

\begin{definitionbox}
型 \(A\) と型 \(B\) の間の関係とは、値 \(a : A\) と \(b : B\) を受け取り、命題を返すものとして読める。
\[
  R : A \to B \to \mathrm{Prop}
\]
同じ型同士の関係 \(R : A \to A \to \mathrm{Prop}\) は、等しさ、大小関係、包含関係、互換性、依存関係などを表せる。
\end{definitionbox}

\begin{leanbox}
Lean では、関係を単なる関数として書ける。返り値が \texttt{Bool} ではなく \texttt{Prop} である点が重要である。\texttt{Prop} にすることで、「この二つが関係している」という主張を定理として証明できる。
\end{leanbox}

\begin{listing}[htbp]
\begin{minted}{lean4}
namespace Ch16

-- A と B の間の関係は、A の値と B の値から命題を作るもの。
def Rel (A B : Type) := A -> B -> Prop

structure DetailLog where
  userId : Nat
  endpoint : Nat
  token : Nat
deriving Repr

structure PublicLog where
  endpoint : Nat
deriving Repr

-- 詳細ログから公開ログへ落とす。
def maskLog (l : DetailLog) : PublicLog :=
  { endpoint := l.endpoint }

-- 公開ログに残すと決めた部分が保存されている、という関係。
def KeepsPublicPart (d : DetailLog) (p : PublicLog) : Prop :=
  p.endpoint = d.endpoint

theorem maskLog_keeps_public (d : DetailLog) :
    KeepsPublicPart d (maskLog d) := by
  rfl
\end{minted}
\caption{詳細ログと公開ログの間の関係}
\label{lst:ch16-relation}
\end{listing}

この定理が示しているのは、\texttt{maskLog} が詳細ログを完全に保存するということではない。\texttt{token} や \texttt{userId} は公開ログから消えている。示しているのは、公開ログに残すと決めた \texttt{endpoint} については、元の詳細ログと対応しているということである。

\begin{warningbox}
情報を落とす処理に対して「元に戻せないから正しくない」と判断すると、抽象化、匿名化、集計、公開用ビューをすべて不正なものとして扱ってしまう。必要なのは、等式ではなく、その処理にふさわしい関係を定義することである。
\end{warningbox}

\section{前順序：反射性と推移性を持つ比較}

次に、同じ型の値を比較する関係を考える。たとえばログの詳細度を、\texttt{public}、\texttt{internal}、\texttt{secret} の三段階で表すとする。このとき、\texttt{public} は \texttt{internal} より情報が少なく、\texttt{internal} は \texttt{secret} より情報が少ない、と読める。

ここで重要なのは、比較が必ずしも「大小」だけを意味しないことである。前順序は、反射性と推移性を満たす関係であればよい。

\begin{definitionbox}
型 \(A\) 上の関係 \(\sqsubseteq\) が前順序であるとは、次の二つを満たすことである。
\[
\begin{aligned}
&\text{反射性:} && a \sqsubseteq a \\
&\text{推移性:} && a \sqsubseteq b \text{ かつ } b \sqsubseteq c \text{ なら } a \sqsubseteq c
\end{aligned}
\]
反射性は「自分自身とは比較できる」ことを意味する。推移性は「安全な比較を二段つなげても安全である」ことを意味する。
\end{definitionbox}

\begin{table}[htbp]
\centering
\caption{前順序として読める実務上の比較}
\label{tab:ch16-preorder-examples}
\begin{tabular}{lll}
\toprule
比較したいもの & \(x \sqsubseteq y\) の読み方 & 実務上の意味 \\
\midrule
ログ詳細度 & \(x\) は \(y\) より情報が少ない & 公開可能性、マスキング \\
仕様 & \(x\) は \(y\) より強い仕様である & 事前条件・事後条件の比較 \\
権限 & \(x\) は \(y\) に含まれる権限である & role / permission 設計 \\
静的解析結果 & \(x\) は \(y\) より精密な近似である & 抽象解釈、警告の安全性 \\
データ形式 & \(x\) は \(y\) へ安全に弱化できる & lossy migration \\
\bottomrule
\end{tabular}
\end{table}

\begin{listing}[htbp]
\begin{minted}{lean4}
-- 前順序の最小構造。
-- le は比較関係、refl は反射性、trans は推移性を表す。
structure PreorderMini (A : Type) where
  le : A -> A -> Prop
  refl : forall a, le a a
  trans : forall {a b c}, le a b -> le b c -> le a c
\end{minted}
\caption{前順序の最小定義}
\label{lst:ch16-preorder-mini}
\end{listing}

本章では、\(x \sqsubseteq y\) を「\(x\) は \(y\) 以下である」と読む。ただし、何が「以下」なのかは場面ごとに決める。ログの例では「情報量が少ない」が以下である。権限の例では「許可が少ない」が以下である。仕様の例では、向きを間違えないよう注意が必要である。一般には、強い仕様は弱い仕様を含意するので、強い仕様を小さい側に置くか大きい側に置くかを、章やプロジェクト内で統一しておく。

\FigurePlaceholder{fig-ch16-preorder-chain.png}{0.88}{ログ詳細度を前順序として見る}{fig:ch16-preorder-chain}

\begin{listing}[htbp]
\begin{minted}{lean4}
inductive Detail where
  | public
  | internal
  | secret
deriving Repr, DecidableEq

-- public <= internal <= secret と読む。
def detailLe : Detail -> Detail -> Prop
  | .public, _ => True
  | .internal, .internal => True
  | .internal, .secret => True
  | .secret, .secret => True
  | _, _ => False

theorem detail_refl : forall d : Detail, detailLe d d := by
  intro d
  cases d <;> simp [detailLe]

theorem detail_trans :
    forall {a b c : Detail}, detailLe a b -> detailLe b c -> detailLe a c := by
  intro a b c hab hbc
  cases a <;> cases b <;> cases c <;> simp [detailLe] at *
\end{minted}
\caption{ログ詳細度の前順序}
\label{lst:ch16-detail-preorder}
\end{listing}

この順序では、\texttt{public} はどのレベル以下でもある。公開情報は、内部情報や秘密情報よりも情報量が少ないからである。一方、\texttt{secret} は \texttt{public} 以下ではない。秘密情報を公開情報として扱うには、抽象化やマスキングが必要になる。

\begin{listing}[htbp]
\begin{minted}{lean4}
theorem no_secret_as_public :
    detailLe Detail.secret Detail.public -> False := by
  intro h
  simpa [detailLe] using h
\end{minted}
\caption{秘密情報をそのまま公開扱いにできないこと}
\label{lst:ch16-no-downgrade}
\end{listing}

\section{仕様の強弱と単調性}

前順序は、データの情報量だけでなく、仕様の強弱にも使える。たとえば、ある入力 \texttt{x} に対する述語 \texttt{p x} と \texttt{q x} を考える。\texttt{p} が成り立つなら必ず \texttt{q} も成り立つとき、\texttt{p} は \texttt{q} より強い仕様であると言える。

\begin{intuitionbox}
強い仕様は、許す値が少ない。弱い仕様は、許す値が多い。たとえば「自然数である」より「偶数の自然数である」のほうが強い。強い仕様を満たす値は、自動的に弱い仕様も満たす。
\end{intuitionbox}

\begin{listing}[htbp]
\begin{minted}{lean4}
-- p が q より強いとは、p を満たせば q も満たすということ。
def Stronger {A : Type} (p q : A -> Prop) : Prop :=
  forall x, p x -> q x

theorem stronger_refl {A : Type} (p : A -> Prop) :
    Stronger p p := by
  intro x hx
  exact hx

theorem stronger_trans {A : Type} {p q r : A -> Prop} :
    Stronger p q -> Stronger q r -> Stronger p r := by
  intro hpq hqr x hpx
  exact hqr x (hpq x hpx)
\end{minted}
\caption{仕様の強弱を含意で表す}
\label{lst:ch16-stronger-spec}
\end{listing}

変換が順序を壊さないことも重要である。この性質を単調性という。ログ詳細度の例では、詳細レベルを公開ビューへ抽象化するとき、元の詳細度の比較が、抽象後の比較としても成り立ってほしい。

\begin{definitionbox}
順序 \(\sqsubseteq_A\) を持つ型 \(A\) から、順序 \(\sqsubseteq_B\) を持つ型 \(B\) への関数 \(f : A \to B\) が単調であるとは、次を満たすことである。
\[
  x \sqsubseteq_A y \quad\Rightarrow\quad f(x) \sqsubseteq_B f(y)
\]
単調性は、「比較できていたものを変換後も比較できる」ことを保証する。
\end{definitionbox}

\begin{listing}[htbp]
\begin{minted}{lean4}
def MonotoneMini {A B : Type}
    (leA : A -> A -> Prop) (leB : B -> B -> Prop)
    (f : A -> B) : Prop :=
  forall {x y}, leA x y -> leB (f x) (f y)
\end{minted}
\caption{単調性の最小定義}
\label{lst:ch16-monotone}
\end{listing}

\section{抽象化と具体化}

ここから、Galois 接続の準備に入る。Galois 接続では、二つの世界を考える。一つは具体的な世界であり、もう一つは抽象的な世界である。

詳細ログの例では、具体的な世界は \texttt{Detail} である。\texttt{public}、\texttt{internal}、\texttt{secret} の三段階を区別する。一方、公開判定のためには、そこまで細かい区別が不要な場合がある。公開ビューでは、「公開情報だけ」か「非公開情報を含みうる」かの二段階だけを見る。

\begin{softwarebox}
抽象化は、実装上は「情報を落とす関数」である。しかし、証明上は「どの情報を落としてよいか」を明示する境界でもある。具体化は、抽象的な判定を満たす具体値の範囲を戻して見る操作として読む。
\end{softwarebox}

\begin{listing}[htbp]
\begin{minted}{lean4}
inductive PublicView where
  | publicOnly
  | mayContainPrivate
deriving Repr, DecidableEq

-- publicOnly <= mayContainPrivate と読む。
def viewLe : PublicView -> PublicView -> Prop
  | .publicOnly, _ => True
  | .mayContainPrivate, .mayContainPrivate => True
  | .mayContainPrivate, .publicOnly => False

-- 抽象化: 詳細レベルを公開判定用の二段階へ落とす。
def erase : Detail -> PublicView
  | .public => .publicOnly
  | .internal => .mayContainPrivate
  | .secret => .mayContainPrivate

-- 具体化: 抽象ビューが許す最大の詳細レベルを見る。
def allow : PublicView -> Detail
  | .publicOnly => .public
  | .mayContainPrivate => .secret
\end{minted}
\caption{詳細度から公開ビューへの抽象化}
\label{lst:ch16-view}
\end{listing}

\texttt{erase} は抽象化である。\texttt{internal} と \texttt{secret} を同じ \texttt{mayContainPrivate} にまとめるので、情報を失う。したがって \texttt{erase} は同型ではない。

一方、\texttt{allow} は具体化である。\texttt{publicOnly} が許す具体レベルの最大値は \texttt{public} であり、\texttt{mayContainPrivate} が許す最大値は \texttt{secret} である、と定めている。

\begin{listing}[htbp]
\begin{minted}{lean4}
theorem view_refl : forall v : PublicView, viewLe v v := by
  intro v
  cases v <;> simp [viewLe]

theorem erase_monotone :
    MonotoneMini detailLe viewLe erase := by
  intro x y h
  cases x <;> cases y <;> simp [detailLe, viewLe, erase] at *

theorem allow_monotone :
    MonotoneMini viewLe detailLe allow := by
  intro x y h
  cases x <;> cases y <;> simp [viewLe, detailLe, allow] at *
\end{minted}
\caption{抽象化が単調であること}
\label{lst:ch16-erase-monotone}
\end{listing}

単調性により、詳細度の順序を抽象化後にも壊さないことがわかる。ただし、単調性だけでは、抽象化と具体化が十分に対応しているとは言えない。そこで Galois 接続を使う。

\FigurePlaceholder{fig-ch16-abstraction-concretization.png}{0.92}{詳細な世界と抽象的な世界を抽象化・具体化で結ぶ}{fig:ch16-abstraction-concretization}

\section{Galois 接続：抽象化と具体化の対応}

Galois 接続は、抽象化 \(\alpha\) と具体化 \(\gamma\) が、順序の比較を通じてぴったり対応していることを表す法則である。

具体的な世界を \(C\)、抽象的な世界を \(A\) とする。それぞれに順序 \(\sqsubseteq_C\)、\(\sqsubseteq_A\) があるとする。抽象化 \(\alpha : C \to A\) と具体化 \(\gamma : A \to C\) が Galois 接続をなすとは、次が成り立つことである。

\[
  \alpha(c) \sqsubseteq_A a
  \quad\Longleftrightarrow\quad
  c \sqsubseteq_C \gamma(a)
\]

左辺は「具体値 \(c\) を抽象化した結果は、抽象値 \(a\) 以下である」と読む。右辺は「具体値 \(c\) は、抽象値 \(a\) を具体化した上限以下である」と読む。この二つが同値であるとき、抽象側での判定と具体側での判定がずれていないと言える。

\begin{definitionbox}
本章で使う Galois 接続は、次の形で覚えればよい。
\[
  \text{抽象化したものが抽象仕様を満たす}
  \quad\Longleftrightarrow\quad
  \text{元の具体値が、その抽象仕様の具体化に収まる}
\]
これは、静的解析、権限設計、公開ログ、仕様弱化で現れる「抽象的に安全なら、具体的にも許された範囲にいる」という対応を表す。
\end{definitionbox}

\begin{listing}[htbp]
\begin{minted}{lean4}
structure GaloisConnectionMini {C A : Type}
    (cLe : C -> C -> Prop) (aLe : A -> A -> Prop)
    (alpha : C -> A) (gamma : A -> C) : Prop where
  law : forall c a, Iff (aLe (alpha c) a) (cLe c (gamma a))

-- erase と allow は Galois 接続をなす。
theorem erase_allow_galois :
    GaloisConnectionMini detailLe viewLe erase allow := by
  constructor
  intro c a
  cases c <;> cases a <;> simp [erase, allow, detailLe, viewLe]
\end{minted}
\caption{Galois接続の最小定義}
\label{lst:ch16-galois-def}
\end{listing}

この証明は、すべての場合を列挙している。\texttt{Detail} は三通り、\texttt{PublicView} は二通りなので、合計六通りである。各ケースで、左辺と右辺が同じ真偽になることを \texttt{simp} で確認している。

\FigurePlaceholder{fig-ch16-galois-law.png}{0.94}{Galois接続の法則：抽象側の比較と具体側の比較が対応する}{fig:ch16-galois-law}

\section{Galois 接続から得られる二つの安全性}

Galois 接続があると、抽象化と具体化の往復について、二つの基本的な性質が得られる。

一つ目は、具体値 \(c\) を抽象化してから具体化すると、元の \(c\) を覆う上限になるという性質である。
\[
  c \sqsubseteq_C \gamma(\alpha(c))
\]
これは「抽象化で情報を落としたあとに戻しても、元の値を含む安全な範囲になる」と読める。

二つ目は、抽象値 \(a\) を具体化してから抽象化すると、元の \(a\) 以下になるという性質である。
\[
  \alpha(\gamma(a)) \sqsubseteq_A a
\]
これは「抽象仕様から許される最大の具体値を選び、それを再び抽象化しても、元の抽象仕様を超えない」と読める。

\begin{listing}[htbp]
\begin{minted}{lean4}
theorem gc_unit {C A : Type} {cLe : C -> C -> Prop} {aLe : A -> A -> Prop}
    {alpha : C -> A} {gamma : A -> C}
    (gc : GaloisConnectionMini cLe aLe alpha gamma)
    (a_refl : forall a, aLe a a) :
    forall c, cLe c (gamma (alpha c)) := by
  intro c
  exact (gc.law c (alpha c)).mp (a_refl (alpha c))

theorem gc_counit {C A : Type} {cLe : C -> C -> Prop} {aLe : A -> A -> Prop}
    {alpha : C -> A} {gamma : A -> C}
    (gc : GaloisConnectionMini cLe aLe alpha gamma)
    (c_refl : forall c, cLe c c) :
    forall a, aLe (alpha (gamma a)) a := by
  intro a
  exact (gc.law (gamma a) a).mpr (c_refl (gamma a))
\end{minted}
\caption{Galois接続から得られる一般的な性質}
\label{lst:ch16-galois-consequences}
\end{listing}

\begin{listing}[htbp]
\begin{minted}{lean4}
theorem detail_covered_by_allowed (d : Detail) :
    detailLe d (allow (erase d)) := by
  exact gc_unit erase_allow_galois view_refl d

theorem allowed_abstraction_is_inside (v : PublicView) :
    viewLe (erase (allow v)) v := by
  exact gc_counit erase_allow_galois detail_refl v
\end{minted}
\caption{公開ログ例に適用した安全性}
\label{lst:ch16-galois-use}
\end{listing}

\texttt{detail\_covered\_by\_allowed} は、詳細レベルを一度抽象化し、その抽象ビューが許す最大の具体レベルへ戻すと、元の詳細レベルを覆うことを言っている。たとえば \texttt{internal} を抽象化すると \texttt{mayContainPrivate} になり、それを具体化すると \texttt{secret} になる。\texttt{internal <= secret} なので、安全に覆われている。

\texttt{allowed\_abstraction\_is\_inside} は、抽象ビューから許される最大の具体レベルを選び、それを再度抽象化しても、元の抽象ビューを超えないことを言っている。これは公開ポリシーの上限を越えないことに対応する。

\section{サンプル：詳細ログから公開ログへの安全な抽象化}

ここまでの道具を、公開ログの設計レビューとして読み替える。

まず、詳細ログには多くの情報が含まれている。公開ログには、その一部だけを残す。この時点で、詳細ログと公開ログは同型ではない。公開ログから詳細ログは復元できない。

次に、公開ログに残す部分については関係で仕様を書く。たとえば、\texttt{endpoint} は保存されているべきである。この性質は \texttt{KeepsPublicPart} として表した。

さらに、ログ全体の安全性は前順序で見る。\texttt{public <= internal <= secret} という順序を決めれば、「公開側に出してよい情報か」「内部に留めるべき情報か」を比較できる。

最後に、詳細な分類から公開判定用の分類へ落とす抽象化 \texttt{erase} と、その抽象分類が許す具体的な最大値 \texttt{allow} を定義する。これらが Galois 接続をなすことを証明すれば、抽象側のポリシー判定と具体側の情報量判定が対応していることを確認できる。

\FigurePlaceholder{fig-ch16-public-log-mask.png}{0.94}{詳細ログから公開ログへ落とすときに証明する性質}{fig:ch16-public-log-mask}

\begin{softwarebox}
設計レビューでは、次のように分けて確認する。
\begin{itemize}
\item 保存すべき項目は、関係として定義する。
\item 情報量や権限の大小は、前順序として定義する。
\item 変換が順序を壊さないことは、単調性として確認する。
\item 抽象化と具体化の対応は、Galois 接続として確認する。
\end{itemize}
これにより、「公開ログは詳細ログと同じではないが、仕様上安全な近似である」と言える。
\end{softwarebox}

\section{実務への読み替え}

Galois 接続は、本章の小さな例だけのものではない。実務では、次のような形で現れる。

\begin{table}[htbp]
\centering
\caption{Galois接続として読める実務例}
\label{tab:ch16-galois-applications}
\begin{tabular}{llll}
\toprule
領域 & 具体側 & 抽象側 & 読み方 \\
\midrule
公開ログ & 詳細ログ & 公開ビュー & マスキング後の安全性 \\
権限設計 & 個別 permission & role & role が許す権限範囲 \\
静的解析 & 実行時状態 & 解析上の近似 & 警告が安全側に倒れていること \\
仕様弱化 & 強い仕様 & 公開・互換仕様 & 弱めた仕様で保証する範囲 \\
データ集約 & 生ログ & 集計値 & 集計が元データの上限・下限を保つこと \\
\bottomrule
\end{tabular}
\end{table}

静的解析の例では、具体側は実際のプログラム状態であり、抽象側は「この変数は非負らしい」「この値は null かもしれない」のような解析結果である。解析は完全に実行時状態を再現しない。むしろ、実行時に起こりうる状態を安全に覆う近似であることが重要である。

権限設計では、具体側を個別 permission の集合、抽象側を role として読むことができる。role を具体化すると、その role が許す permission の集合になる。個別権限が role の範囲内に収まっているかどうかは、前順序や包含関係で表せる。

lossy migration では、完全な round-trip は期待できない。第23章では、この章の考え方を使って、復元ではなく仕様弱化や正規化後の一致として安全性を表す。

\section{よくある誤解}

\begin{warningbox}
\textbf{誤解1：情報を落とす処理はすべて悪い。}\par
情報を落とすこと自体は悪くない。公開、匿名化、集計、抽象化では、情報を落とすことが目的である。問題は、どの情報を落としてよく、どの性質を残すべきかを明示していないことである。
\end{warningbox}

\begin{warningbox}
\textbf{誤解2：Galois 接続は難しい圏論の話なので実務には関係ない。}\par
本章で使った形では、Galois 接続は「抽象側での判定」と「具体側での判定」が一致することを表す道具である。静的解析、権限、公開ログのように、具体的なものを抽象的な分類で扱う場面では自然に現れる。
\end{warningbox}

\begin{warningbox}
\textbf{誤解3：順序の向きはどちらでもよい。}\par
数学的には向きを反対にしても理論は作れる。しかし、プロジェクト内で向きが揺れると、定理名やレビューコメントが読みにくくなる。本章では「情報が少ないほうを小さい」とした。仕様の強弱や権限では、どちらを小さい側に置くかを最初に決める必要がある。
\end{warningbox}

\section{本章のまとめ}

関係は、二つの値が何らかの意味で対応していることを表す命題である。

前順序は、反射性と推移性を持つ比較であり、情報量、権限、仕様の強弱、抽象度を表せる。

単調性は、変換が比較関係を壊さないことを表す。

Galois 接続は、抽象化と具体化が順序を通じて対応していることを表す。

情報を落とす変換は、同型ではなく、安全な近似として仕様化すべき場合がある。

本章では、等式から一歩進み、順序や近似で安全性を扱う方法を学んだ。次章では、複数のデータが共通キーや整合条件で合流する場面を扱う。そこでは、Pullback を「制約を満たすペア」として読み、DB join や複数APIレスポンスの統合へ接続する。
